% Chapter 24: The Kakeya Conjecture
\let\LocalMacro\relax

\chapter{The Kakeya Conjecture}

\section{The Problem}
In 1917, Soichi Kakeya posed a simple question in geometry \cite{kakeya1917}: \emph{What is the minimum area of a planar set in which a unit line segment can be turned 360 degrees (a needle set)?}

In 1919, Abram Besicovitch shocked the mathematical community by proving that such a set—now called a \emph{Besicovitch set} or \emph{Kakeya set}—can have Lebesgue measure **zero** \cite{besicovitch1919}:

\begin{definition}[Kakeya Set]
A subset $K \subset \mathbb{R}^n$ is a \emph{Kakeya set} if it contains a unit line segment in every direction. That is, for every unit vector $e \in S^{n-1}$, there exists $x \in \mathbb{R}^n$ such that $\{x + t e : t \in [0, 1]\} \subset K$.
\end{definition}

Since Kakeya sets in $\mathbb{R}^n$ can have measure zero, mathematicians shifted their focus to the \emph{dimension} of these sets. This led to the Kakeya Conjecture:

\begin{conjecture}[Kakeya Conjecture]
Any Kakeya set in $\mathbb{R}^n$ must have Hausdorff dimension and Minkowski dimension equal to $n$.
\end{conjecture}

Despite its geometric statement, the Kakeya conjecture is deeply connected to harmonic analysis, specifically to wave equations, Fourier multipliers, and the restriction conjecture.

\begin{historicalbox}
The Kakeya conjecture is one of the most famous unsolved problems in analysis. While it is settled for $n=2$, it remains open in higher dimensions ($n \geq 3$), representing a key barrier to understanding multi-directional wave propagation.
\end{historicalbox}

\section{Mathematical Theory}


\subsection{Prerequisites}
This chapter assumes familiarity with:
\begin{itemize}
\item Basic real analysis (Lebesgue measure, open/closed sets, limits)
\item Geometric measure theory (Hausdorff dimension)
\item Finite fields and polynomials (roots, finite vector spaces)
\end{itemize}

\subsection{Besicovitch's Construction}
Besicovitch's proof uses a geometric decomposition called "panning." By dividing a triangle into smaller triangles, shifting them, and overlapping them, one can concentrate the area into arbitrarily small regions while preserving segments in all directions.

\subsection{Finite Field Analogue}
In 1999, Thomas Wolff proposed studying the Kakeya conjecture in the simplified setting of vector spaces over finite fields $\mathbb{F}_q^n$ to bypass real-variable complications \cite{wolff1999}.
\begin{definition}[Finite Field Kakeya Set]
A subset $K \subset \mathbb{F}_q^n$ is a Kakeya set if for every direction $y \in \mathbb{F}_q^n \setminus \{0\}$, there exists $x \in \mathbb{F}_q^n$ such that $\{x + t y : t \in \mathbb{F}_q\} \subset K$.
\end{definition}
Wolff conjectured that any Kakeya set in $\mathbb{F}_q^n$ satisfies $|K| \geq C_n q^n$.

\section{The Discovery}

The Kakeya conjecture has seen two major breakthroughs in recent years:

\subsection{Dvir's Finite Field Proof (2008)}
In 2008, Zeev Dvir proved Wolff's finite field conjecture in a paper that was only four pages long, using the \emph{polynomial method} \cite{dvir2009}:

\begin{theorem}[Dvir, 2008 \cite{dvir2009}]
Any Kakeya set $K \subset \mathbb{F}_q^n$ satisfies:
\begin{equation}
|K| \geq \frac{1}{n!} q^n
\end{equation}
\end{theorem}
The proof is elegant: if $|K| < \binom{n+q-1}{n}$, there exists a non-zero polynomial of degree $<q$ vanishing on $K$. By complete induction along line directions, the polynomial's highest degree homogeneous part must vanish on all of $\mathbb{F}_q^n$, which is impossible.

\subsection{The 3D Real Breakthrough (2023--2025)}
In 2023 and 2024, Hong Wang (NYU) and Joshua Zahl (UBC) made the first major progress on the real conjecture in three dimensions in decades \cite{wangzahl2024}:

\begin{theorem}[Wang--Zahl, 2024]
Any Kakeya set in $\mathbb{R}^3$ has Hausdorff dimension at least $2.5 + \epsilon$ for a positive constant $\epsilon > 0$.
\end{theorem}

Their proof combined Dvir's polynomial method with \emph{polynomial partitioning} (introduced by Guth and Katz in 2010) and linear programming bounds on line incidences, breaking the long-standing "2.5 dimension barrier" in $\mathbb{R}^3$.

\begin{highlightbox}
The Kakeya conjecture progress in 3D by Wang and Zahl represents a major triumph of polynomial techniques in real harmonic analysis, showing how algebraic structures can control geometric sets in $\mathbb{R}^3$.
\end{highlightbox}

\section{Impact}
\begin{itemize}
\item \textbf{Harmonic Analysis}: Directly bounds the restriction operator for the Fourier transform, influencing local smoothing estimates for wave equations.
\item \textbf{Additive Combinatorics}: The finite field polynomial method has been adapted to prove the joints conjecture and other incidence geometry problems.
\item \textbf{Partial Differential Equations}: Provides bounds on the size of singularity sets for dispersive wave propagation.
\end{itemize}

\section{Exercises}

\begin{exercise}[Basic]
Show that a Kakeya set in $\mathbb{R}^2$ containing segments in all directions can be constructed with measure 0.
\end{exercise}

\begin{exercise}[Intermediate]
Write down the details of Zeev Dvir's finite field Kakeya set proof. Show why the degree of the vanishing polynomial must be less than $q$.
\end{exercise}

\begin{exercise}[Challenge]
Explain why Dvir's algebraic proof fails in the setting of real numbers $\mathbb{R}^n$. What is the main difference between lines in $\mathbb{F}_q^n$ and $\mathbb{R}^n$?
\end{exercise}

\begin{exercise}[Computational]
Write a Python/SageMath script to explore a concept from this chapter. See the appendix for starter code.
\end{exercise}

\section{Timeline}

\begin{tabularx}{\linewidth}{lX}
\toprule
\textbf{Year} & \textbf{Event} \\ \midrule
1917 & Kakeya poses the needle turning problem \cite{kakeya1917} \\ 
1919 & Besicovitch constructs a Kakeya set of measure zero \cite{besicovitch1919} \\ 
1999 & Wolff formulates the finite field Kakeya conjecture \cite{wolff1999} \\ 
2008 & Dvir proves the finite field Kakeya conjecture using the polynomial method \cite{dvir2009} \\ 
2023--2025 & Wang and Zahl break the 2.5-dimension barrier in $\mathbb{R}^3$ \cite{wangzahl2024} \\
2026 & Hong Wang awarded the Fields Medal for contributions to harmonic analysis and the Kakeya conjecture \\ \bottomrule
\end{tabularx}

\section{Citations}

\begin{enumerate}
\item Kakeya, S. (1917). ``Some problems on maximum and minimum regarding ovals.'' Tohoku Mathematical Journal.
\item Besicovitch, A. S. (1919). ``Sur deux questions de geometrie.'' Bulletin de l'Academie des Sciences de Russie.
\item Wolff, T. (1999). ``Recent work on restriction and Kakeya conjectures.'' IAS Lectures.
\item Dvir, Z. (2009). ``On the size of Kakeya sets in finite fields.'' Journal of the AMS, 23(4), 1107--1119.
\item Wang, H., \& Zahl, J. (2025). ``The Kakeya conjecture in three dimensions.'' arXiv:2502.17655.
\item International Mathematical Union. (2026). ``Fields Medal citation: Hong Wang.'' ICM 2026, Philadelphia.
\end{enumerate}
